% 原书物理页=印刷页2--5；1.1 全节（1.1.1--1.1.15；其中 1.1.14--1.1.15 位于第 5 页）.
\chapter{一般概念}\label{chap:fg-1}
\section{群、子群、同态与商}\label{sec:fg-1-1}

\begin{definition}\label{def:fg-1-1-1}
一个\newterm{fg-group}{群}{group}是一个集合，其上带有一个\newterm{fg-internal-law}{内部运算}{internal law of composition}，即一个映射
\[
\begin{array}{rcl}
G\times G & \longrightarrow & G\\
(x,y) & \longmapsto & xy
\end{array},
\]
此外还满足下列性质：
\begin{itemize}[label=--]
\item 结合律：对一切 $x,y,z\in G$，有 $(xy)z=x(yz)$，因此可把它记作 $xyz$.
\item 存在一个\newterm{fg-neutral-element}{单位元}{neutral element} $e$，使得 $\forall x\in G,\ ex=xe=x$.
\item 存在\newterm{fg-inverse}{逆元}{inverse}：$\forall x\in G,\ \exists y\in G,\ xy=yx=e$. 立即验证可知这样的 $y$ 是唯一的，我们称之为``$x$ 的逆元''，并记作 $x^{-1}$.
\end{itemize}

该运算也可以用不同的符号来记，如 $(x,y)\mapsto x\circ y$、$x*y$ 或 $x\cdot y$. 当群是交换的时候，人们乐意把它记作 $x+y$（此时 $x$ 的逆元记作 $-x$，也称为``$x$ 的负元''）. 关于此，提醒以下两点：
\begin{itemize}[label=--]
\item $G$ 的两个元素 $x,y$ 交换，是指 $xy=yx$；
\item 群 $G$ 是 \newterm{fg-abelian-group}{Abel 群}{abelian group}（或交换群），如果这对一切 $x,y\in G$ 都成立.
\end{itemize}
当 $G$ 是有限集合时，它的基数（即 $G$ 的元素个数）称为 $G$ 的\newterm{fg-order}{阶}{order}，记作 $|G|$.
\end{definition}

\begin{definition}\label{def:fg-1-1-2}
$G$ 的一个\newterm{fg-subgroup}{子群}{subgroup} $H$（记号：$H<G$）是指 $G$ 的一个包含 $e$ 的子集，它在 $G$ 的群运算下封闭且对取逆封闭；等价地，即满足对一切 $x,y\in H$ 有 $xy^{-1}\in H$.
\end{definition}

容易验证：一些子群的交仍是子群. 于是设 $S\subset G$ 为 $G$ 的任一子集. $G$ 中所有包含 $S$ 的子群之交 $\langle S\rangle$ 称为由 $S$ 生成的子群. 它是 $G$ 中包含 $S$ 的``最小''子群. 当 $S=\{x\}$ 是单元素集时，把由 $x$ 生成的子群简单地记作 $\langle x\rangle:=\langle\{x\}\rangle$. 基数 $|\langle x\rangle|$ 称为 $x$ 的阶，有时记作 $\operatorname{ord}(x)$.

若 $G$ 由单独一个元素生成，则称 $G$ 是\newterm{fg-monogeneous}{由一个元素生成的}{monogeneous}；若 $G$ 由一个元素生成且具有有限阶，则称 $G$ 是\newterm{fg-cyclic}{循环的}{cyclic}.

\begin{definition}\label{def:fg-1-1-3}
从群 $G$ 到群 $G'$ 的一个\newterm{fg-homomorphism}{（群）同态}{(group) homomorphism}是一个映射 $f:G\to G'$，使得对一切 $(x,y)\in G\times G$ 有 $f(xy)=f(x)f(y)$.
\end{definition}

此时对一切 $x$ 有 $f(e)=e'$ 与 $f(x^{-1})=f(x)^{-1}$. 我们也常用 morphisme 一词代替 homomorphisme（二者中文均译作``同态''）. 若 $G=G'$，则称 $f$ 为\newterm{fg-endomorphism}{自同态}{endomorphism}；若存在从 $G'$ 到 $G$ 的一个同态 $g$ 使得 $f\circ g=\mathrm{id}_{G'}$、$g\circ f=\mathrm{id}_G$，则称 $f$ 为\newterm{fg-isomorphism}{同构}{isomorphism}；若 $f$ 同时是自同态与同构，则称 $f$ 为\newterm{fg-automorphism}{自同构}{automorphism}.

立即验证：$G$ 的子群 $H$ 的\newterm{fg-image}{像}{image} $f(H)$ 是 $G'$ 的子群. 特别地，$f(G)$（我们也把它记作 $\operatorname{Im}(f)$）就是这种情形.

同样，$G'$ 的子群 $H'$ 的逆像（或原像）$f^{-1}(H')$ 是 $G$ 的子群. 特别地，$f^{-1}(\{e'\})$ 就是这种情形，我们也把它记作 $\operatorname{Ker}(f)$，并称之为 $f$ 的\newterm{fg-kernel}{核}{kernel}.

\begin{numpara}\label{np:fg-1-1-4}
设 $A$ 与 $B$ 为 $G$ 的两个子集. 我们把 $G$ 中形如 $ab$（$a\in A$，$b\in B$）的元素（相应地，形如 $a^{-1}$，$a\in A$）构成的子集记作 $AB$（相应地，$A^{-1}$）. 换言之，我们有
\[
AB=\{x\in G,\ \exists(a,b)\in A\times B,\ x=ab\}\quad\text{和}\quad A^{-1}=\{a^{-1},\ a\in A\}.
\]
例如，$A$ 是 $G$ 的子群当且仅当 $A$ 非空且 $A.A^{-1}=A$. 若 $A$ 与 $B$ 是 $G$ 的子群，则 $AB$ 是 $G$ 的子群当且仅当 $AB=BA$.
\end{numpara}

\begin{numpara}\label{np:fg-1-1-5}
模 $H$ 的\newterm{fg-left-coset}{左陪集}{left coset}是指 $G$ 的形如
\[
C=xH=\{xh;\ h\in H\}
\]
的子集 $C$，其中 $x$ 是 $G$ 的一个元素. 此时对任意元素 $y\in C$ 也有 $C=yH$. $C$ 的元素也称为 $C$ 在 $G$ 中的``代表元''. 模 $H$ 的左陪集全体记作 $G/H$. 它是 $G$ 的子集全体构成的集合 $\mathcal{P}(G)$ 的一个子集，且我们有一个满的``典范''映射
\[
G\xrightarrow{\ \pi\ }G/H,\quad x\mapsto xH,
\]
它的纤维恰好就是模 $H$ 的左陪集.

设 $G$ 有限. $G/H$ 的基数记作 $[G:H]$，称为 $H$ 在 $G$ 中的\newterm{fg-index}{指数}{index}. 我们同样定义右陪集 $Hx$ 的集合 $H\setminus G$；它同样有 $[G:H]$ 个元素（考虑映射 $g\mapsto g^{-1}$）.

当 $G$ 有限时，模 $H$ 的所有左陪集都具有相同数目的元素，即 $|H|$. 由于两个不同的陪集不相交，我们便得到关系式：
\end{numpara}

\begin{proposition}\label{prop:fg-1-1-6}
$|G|=|H|\times[G:H]$.
\end{proposition}

于是，若 $G$ 有限，则 $|H|$ 整除 $|G|$（Lagrange 定理）；特别地，$G$ 的任一元素的阶整除 $G$ 的阶.

\begin{definition}\label{def:fg-1-1-7}
称 $H$ 是 $G$ 的\newterm{fg-normal-subgroup}{正规子群}{normal subgroup}，如果对一切 $x\in G$ 有 $xH=Hx$. 这等价于 $\forall x\in G,\ xHx^{-1}=H$，或者更明确地：
\[
\forall x\in X,\ \forall h\in H,\ xhx^{-1}\in H.\footnote{译注：原文此处作 $x\in X$，疑为 $x\in G$ 之误.}
\]
我们常记 $H\triangleleft G$，以强调子群 $H<G$ 是正规的.
\end{definition}

\begin{proposition}\label{prop:fg-1-1-8}
设 $H$ 正规. 则在集合 $G/H$ 上存在唯一的群结构，使得典范满射 $\pi:G\to G/H,\ x\mapsto xH$ 是一个群同态.
\end{proposition}

在该命题中得到的群 $G/H$ 称为 $G$ 对 $H$ 的\newterm{fg-quotient-group}{商群}{quotient group}. 它还由下面的性质刻画：

\begin{proposition}\label{prop:fg-1-1-9}
对任意同态 $G\xrightarrow{\ f\ }G'$，若 $f(H)=\{e'\}$，则存在唯一的同态 $\bar f:G/H\to G'$，使得 $f$ 分解为
\[
f:\ G\xrightarrow{\ \pi\ }G/H\xrightarrow{\ \bar f\ }G'.
\]
\end{proposition}

将此应用于 $f$ 的核，便得到同态的典范分解：

\begin{proposition}\label{prop:fg-1-1-10}
设 $f:G\to G'$ 是一个群同态. 则 $\operatorname{Ker}(f)$ 是 $G$ 的正规子群，且前一命题中的同态 $\bar f$ 应用于 $H=\operatorname{Ker}(f)$ 时诱导出一个同构
\[
\bar f:\ G/\operatorname{Ker}(f)\xrightarrow{\ \sim\ }\operatorname{Im}(f),
\]
它嵌入 $f$ 的下列分解之中：
\[
f:\ G\xrightarrow{\ \pi\ }G/\operatorname{Ker}(f)\xrightarrow{\ \bar f\ }\operatorname{Im}(f)\xrightarrow{\ i\ }G',
\]
其中 $\pi$ 是典范投影，$i$ 是 $\operatorname{Im}(f)$ 到 $G$ 中的包含映射\footnote{译注：原文此处作``dans $G$''（到 $G$ 中）；按 $f:G\to G'$ 与 $\operatorname{Im}(f)\subset G'$，应为 $G'$，疑为原书笔误，今按原文照录.}.
\end{proposition}

特别地，对有限群我们有等式 $|G|=|\operatorname{Ker}(f)|\cdot|\operatorname{Im}(f)|$. 最后，下面两个结论将在习题课中证明.

\begin{proposition}\label{prop:fg-1-1-11}
设 $G$ 是群，$H$ 是 $G$ 的一个正规子群，$K$ 是 $G$ 的一个子群. 那么，$KH=HK$ 是 $G$ 的子群，$H$ 是 $KH$ 的正规子群，$H\cap K$ 是 $K$ 的正规子群，且映射 $x\mapsto\pi(x)=xH$ 诱导出同构 $K/K\cap H\simeq KH/H$.
\end{proposition}

\begin{proposition}\label{prop:fg-1-1-12}
设 $G$ 是群，$H$ 是 $G$ 的一个正规子群. 映射 $\pi:G\to G/H$ 在``$G$ 中包含 $H$ 的子群 $A$ 全体''与``$G/H$ 中形如 $A/H$ 的子群全体''之间建立了一个双射，对该双射有 $A\triangleleft G\Leftrightarrow A/H\triangleleft G/H$；此时，映射 $xH\mapsto xA$ 诱导出同构 $(G/H)/(A/H)\simeq G/A$.
\end{proposition}

\begin{numpara}[直积与直和]\label{np:fg-1-1-13}
设 $G,G'$ 为两个群. 我们在 $G\times G'$ 上定义群结构：$(x,x')(y,y'):=(xy,x'y')$. 更一般地，设 $I$ 是一个集合，$\{G_i;\ i\in I\}$ 是一族群. 我们在诸集合 $G_i$ 的乘积上定义一个群结构，称为\newterm{fg-direct-product}{（直）积}{(direct) product}，记作 $\prod_{i\in I}G_i)$\footnote{译注：原书此处作 $\prod_{i\in I}G_i)$，多印了一个右括号，今照录.}，其运算为 $(x_i)_{i\in I}\cdot(y_i)_{i\in I}=(x_iy_i)_{i\in I}$. 若所有的 $G_i$ 都等于同一个 $G$，则该群记作 $G^I$（当 $\mathrm{card}(I)=n$ 有限时记作 $G^n$），它等同于从 $I$ 到 $G$ 的映射全体构成的群. 当诸 $G_i$ 为 Abel 群时，余积这一对偶概念给出\newterm{fg-direct-sum}{直和}{direct sum}的概念：它是 $\prod_{i\in I}G_i$\footnote{译注：原文此处在``de''（法文，即``的''）之后、$\prod_{i\in I}G_i$ 之前衍一逗号，中文行文无法对应照录，特此说明.}的子群 $\bigoplus_{i\in I}G_i$，由那些除有限多个指标 $i\in I$ 外满足 $x_i=e_i$ 的 $(x_i)_{i\in I}$ 组成.
\end{numpara}

\begin{numpara}[内自同构与外自同构]\label{np:fg-1-1-14}
设 $G$ 是一个群. $G$ 的自同构全体在映射的复合运算下构成一个群，记作 $\operatorname{Aut}(G)$. 对任意 $h\in G$，映射
\[
\operatorname{Int}(h):\ G\to G;\quad g\mapsto\operatorname{Int}(h)(g):=hgh^{-1}
\]
是 $G$ 的一个自同构，称为\newterm{fg-inner-automorphism}{内自同构}{inner automorphism}；我们也称它为``用 $h$ 作\newterm{fg-conjugation}{共轭}{conjugation}''.
映射 $\operatorname{Int}:G\to\operatorname{Aut}(G):\ h\mapsto\operatorname{Int}(h)$ 是一个群同态，其核
\[
Z(G):=\{h\in G,\ \forall g\in G,\ hg=gh\}
\]
称为 $G$ 的\newterm{fg-center}{中心}{center}. $\operatorname{Int}$ 的像 $\operatorname{Int}(G)\simeq G/Z(G)$ 是 $\operatorname{Aut}(G)$ 的正规子群，因为对任意 $h\in G$ 与 $f\in\operatorname{Aut}(G)$ 有 $f\circ\operatorname{Int}(h)\circ f^{-1}=\operatorname{Int}(f(h))$. 商群 $\operatorname{Out}(G):=\operatorname{Aut}(G)/\operatorname{Int}(G)$ 的元素称为 $G$ 的\newterm{fg-outer-automorphism}{外自同构}{outer automorphism}.

于是，$G$ 的子群 $H$ 是正规的，如果对所有 $f\in\operatorname{Int}(G)$ 都有 $f(H)=H$. 若对所有 $f\in\operatorname{Aut}(G)$ 都有 $f(H)=H$，则称 $H$ 是\newterm{fg-characteristic-subgroup}{特征子群}{characteristic subgroup}. 例如，$G$ 的中心是一个特征子群.
\end{numpara}

\begin{definition}\label{def:fg-1-1-15}
称群 $G$ 是\newterm{fg-simple-group}{单群}{simple group}，如果它没有真正规子群（即异于 $\{1\}$ 与 $G$ 的正规子群）.
\end{definition}
